\documentclass[9pt]{extarticle}

\usepackage[margin=0.45in]{geometry}
\usepackage{array}
\usepackage{booktabs}
\usepackage{tabularx}
\usepackage{xcolor}
\usepackage{listings}
\usepackage{amsmath}
\usepackage{enumitem}
\usepackage{titlesec}
\usepackage[T1]{fontenc}
\usepackage{lmodern}

\setlength{\parindent}{0pt}
\setlength{\parskip}{3pt}
\setlength{\tabcolsep}{4pt}
\renewcommand{\arraystretch}{1.05}
\pagestyle{plain}

\titleformat{\section}{\normalsize\bfseries}{}{0pt}{}
\titlespacing*{\section}{0pt}{4pt}{2pt}

\lstdefinestyle{code}{
  basicstyle=\ttfamily\footnotesize,
  columns=fullflexible,
  keepspaces=true,
  showstringspaces=false,
  breaklines=true,
  backgroundcolor=\color{gray!8},
  frame=single,
  rulecolor=\color{gray!35},
  framerule=0.3pt,
  xleftmargin=2pt,
  xrightmargin=2pt,
  aboveskip=2pt,
  belowskip=2pt
}

\newcommand{\py}[1]{\lstinline[style=code]!#1!}
\newcommand{\matlab}[1]{\lstinline[style=code]!#1!}

\begin{document}

\begin{center}
{\LARGE \textbf{Python/NumPy/Matplotlib to MATLAB Quick Reference}}\\[3pt]
{\small A compact syntax guide for students moving from Python to MATLAB}
\end{center}

\vspace{-4pt}
\textbf{Main idea:} NumPy arrays and MATLAB matrices are very similar, but indexing differs: Python starts at \textbf{0}, MATLAB starts at \textbf{1}. Python uses functions such as \py{np.sin(x)}, whereas MATLAB often uses nearly the same mathematical name, e.g. \matlab{sin(x)}.

\section*{1. Variables, vectors, and matrices}

\begin{tabularx}{\textwidth}{>{\raggedright\arraybackslash}p{0.28\textwidth} X X}
\toprule
\textbf{Operation} & \textbf{Python / NumPy} & \textbf{MATLAB} \\
\midrule
Scalar & \py{x = 3.0} & \matlab{x = 3.0;} \\
Row vector & \py{x = np.array([1,2,3])} & \matlab{x = [1 2 3];} \\
Column vector & \py{x = np.array([[1],[2],[3]])} & \matlab{x = [1; 2; 3];} \\
Matrix & \py{A = np.array([[1,2],[3,4]])} & \matlab{A = [1 2; 3 4];} \\
Zeros & \py{np.zeros((3,4))} & \matlab{zeros(3,4)} \\
Ones & \py{np.ones((3,4))} & \matlab{ones(3,4)} \\
Identity matrix & \py{np.eye(4)} & \matlab{eye(4)} \\
Sequence with step & \py{np.arange(0,10,2)} & \matlab{0:2:8} \\
Equally spaced values & \py{np.linspace(0,1,6)} & \matlab{linspace(0,1,6)} \\
Matrix size & \py{A.shape} & \matlab{size(A)} \\
Number of elements & \py{A.size} & \matlab{numel(A)} \\
Transpose & \py{A.T} & \matlab{A.'} \\
\bottomrule
\end{tabularx}

\section*{2. Indexing and slicing}

\begin{tabularx}{\textwidth}{>{\raggedright\arraybackslash}p{0.28\textwidth} X X}
\toprule
\textbf{Operation} & \textbf{Python / NumPy} & \textbf{MATLAB} \\
\midrule
First element & \py{x[0]} & \matlab{x(1)} \\
Third element & \py{x[2]} & \matlab{x(3)} \\
First row & \py{A[0,:]} & \matlab{A(1,:)} \\
Second column & \py{A[:,1]} & \matlab{A(:,2)} \\
Rows 2--4 & \py{A[1:4,:]} & \matlab{A(2:4,:)} \\
Last element & \py{x[-1]} & \matlab{x(end)} \\
Every second element & \py{x[::2]} & \matlab{x(1:2:end)} \\
Elements larger than 0 & \py{x[x > 0]} & \matlab{x(x > 0)} \\
\bottomrule
\end{tabularx}

\textbf{Important:} Python slices exclude the final index. For example, \py{x[1:4]} contains indices 1, 2, and 3. MATLAB ranges include the final index, so \matlab{x(2:4)} contains elements 2, 3, and 4.

\section*{3. Basic mathematical operations}

\begin{tabularx}{\textwidth}{>{\raggedright\arraybackslash}p{0.28\textwidth} X X}
\toprule
\textbf{Operation} & \textbf{Python / NumPy} & \textbf{MATLAB} \\
\midrule
Addition & \py{A + B} & \matlab{A + B} \\
Element-by-element multiply & \py{A * B} & \matlab{A .* B} \\
Matrix multiplication & \py{A @ B} & \matlab{A * B} \\
Element-by-element division & \py{A / B} & \matlab{A ./ B} \\
Element-by-element power & \py{A**2} & \matlab{A.^2} \\
Square root & \py{np.sqrt(x)} & \matlab{sqrt(x)} \\
Exponential & \py{np.exp(x)} & \matlab{exp(x)} \\
Sine / cosine & \py{np.sin(x), np.cos(x)} & \matlab{sin(x), cos(x)} \\
Absolute value & \py{np.abs(x)} & \matlab{abs(x)} \\
Sum & \py{np.sum(x)} & \matlab{sum(x)} \\
Mean & \py{np.mean(x)} & \matlab{mean(x)} \\
Maximum & \py{np.max(x)} & \matlab{max(x)} \\
Matrix inverse & \py{np.linalg.inv(A)} & \matlab{inv(A)} \\
Solve $Ax=b$ & \py{x = np.linalg.solve(A,b)} & \matlab{x = A\\b;} \\
\bottomrule
\end{tabularx}

\section*{4. Functions, loops, and conditions}

\begin{tabularx}{\textwidth}{>{\raggedright\arraybackslash}p{0.22\textwidth} X X}
\toprule
\textbf{Operation} & \textbf{Python} & \textbf{MATLAB} \\
\midrule
Simple function &
\texttt{def f(x):}\\[-1pt]
&
\texttt{function y = f(x)}\\[-1pt]
\texttt{\ \ \ \ return x**2 + 1}
&
\texttt{\ \ \ \ y = x.\^{}2 + 1;}\\[-1pt]
&&\texttt{end} \\
\addlinespace
For loop &
\texttt{for i in range(5):}\\[-1pt]
&
\texttt{for i = 1:5}\\[-1pt]
\texttt{\ \ \ \ print(i)}
&
\texttt{\ \ \ \ disp(i)}\\[-1pt]
&&\texttt{end} \\
\addlinespace
If statement &
\texttt{if x > 0:}\\[-1pt]
&
\texttt{if x > 0}\\[-1pt]
\texttt{\ \ \ \ y = 1} & \texttt{\ \ \ \ y = 1;}\\[-1pt]
\texttt{else:} & \texttt{else}\\[-1pt]
\texttt{\ \ \ \ y = -1} & \texttt{\ \ \ \ y = -1;}\\[-1pt]
&\texttt{end}\\
\bottomrule
\end{tabularx}

\section*{5. Plotting: Matplotlib and MATLAB}

\begin{tabularx}{\textwidth}{>{\raggedright\arraybackslash}p{0.24\textwidth} X X}
\toprule
\textbf{Operation} & \textbf{Python / Matplotlib} & \textbf{MATLAB} \\
\midrule
Basic line plot & \py{plt.plot(x,y); plt.show()} & \matlab{plot(x,y)} \\
Line with symbols & \py{plt.plot(x,y,'o-')} & \matlab{plot(x,y,'o-')} \\
Axis labels & \py{plt.xlabel('Time'); plt.ylabel('Amplitude')} & \matlab{xlabel('Time'); ylabel('Amplitude')} \\
Title & \py{plt.title('Signal')} & \matlab{title('Signal')} \\
Grid & \py{plt.grid(True)} & \matlab{grid on} \\
Axis limits & \py{plt.xlim(0,10)} & \matlab{xlim([0 10])} \\
Two curves & \py{plt.plot(x,y1); plt.plot(x,y2)} & \matlab{plot(x,y1); hold on; plot(x,y2)} \\
Legend & \py{plt.legend(['A','B'])} & \matlab{legend('A','B')} \\
Image / matrix & \py{plt.imshow(A, aspect='auto')} & \matlab{imagesc(A)} \\
Colorbar & \py{plt.colorbar()} & \matlab{colorbar} \\
\bottomrule
\end{tabularx}

\section*{6. A small side-by-side example}

\begin{minipage}[t]{0.48\textwidth}
\textbf{Python / NumPy / Matplotlib}
\begin{lstlisting}[style=code]
import numpy as np
import matplotlib.pyplot as plt

x = np.linspace(0, 2*np.pi, 200)
y = np.sin(x)

A = np.array([[1,2],[3,4]])
b = np.array([1,2])
u = np.linalg.solve(A,b)

plt.plot(x,y)
plt.xlabel('x')
plt.ylabel('sin(x)')
plt.grid(True)
plt.show()
\end{lstlisting}
\end{minipage}
\hfill
\begin{minipage}[t]{0.48\textwidth}
\textbf{MATLAB}
\begin{lstlisting}[style=code]
x = linspace(0,2*pi,200);
y = sin(x);

A = [1 2; 3 4];
b = [1; 2];
u = A\b;

plot(x,y)
xlabel('x')
ylabel('sin(x)')
grid on
\end{lstlisting}
\end{minipage}

\vspace{4pt}
\textbf{MATLAB reminders:}
\begin{itemize}[leftmargin=1.2em,itemsep=0pt,topsep=2pt]
    \item A semicolon suppresses output: \matlab{x = 3;} versus \matlab{x = 3}.
    \item Use \matlab{.*}, \matlab{./}, and \matlab{.^} for element-by-element operations.
    \item Use \matlab{*} for matrix multiplication.
    \item MATLAB indexing begins at 1; Python indexing begins at 0.
\end{itemize}

\end{document}
